% Pursuing Stacks (cotes 134-2 to 134-8): archive cote 134-2, site page 39 (PDF page 40).
% appendix (letters to Breen). Draft compared: ps1.tex:1193-1197.
% AI-transcribed by Claude Opus 5.5 (claude-opus-5-5, adaptive thinking, effort high), Message Batches API, batch ps1 (msgbatch_01Rv5hCaNTtUt4JHx1yqoHPL), 2026-11-12; not yet checked by eye.
\scanpage{134-2/39}
\mspage{24}

Tu me demandes des exemples "convainquants" de 2-catégories de Picard.
Voici quelques \deleted{uns}\inserted{exemples,} en vrac \hand{(je ne sais s'ils seront convainquants !)}:

\margin{$\times$(22)} 1) Si L est un lien\inserted{(22)} de centre Z sur le topos X, les \uline{gerbes liées par Z}
forment une 2-catégorie de Picard stricte, représentée par le complexe
$R\hand{\Gamma}_X(Z)$ tronqué en degré 2, dont les objets de cohomologie non triviaux
sont \unsure{Hes|les} $H^i(X,Z)$, $0 \hand{\leqslant} i \hand{\leqslant} 2$. Les \over{\unsure{g}}{g}erbes liées par L forment un \uline{pseudo-2-}
\uline{torseur} sous la gerbe précédente, qui est un 2-torseur\hand{(}i.e. non vide\hand{)} sss
une certaine obstruction dans $H^3(X,Z)$ est nulle. Pour comprendre \unsure{cette}
classe de notre point de vue, il y a lieu de passer aux 2-champs \unsure{corres-|derres-}%
pondants: le 2-champ de\xout{s} Picard \inserted[typed]{strict,} des Z-gerbes\xout{\unsure{liées par Z et le 2-champ des}}
sur des objets variables de X, et le 2-champ des L-gerbes sur des objets
variables. Ce dernier est bel et bien un 2-torseur sous le champ précédent,
or la classification de ces \xout{deux} 2-torseurs (à 2-équivalence près) se
fait par le $H^3(X,Z)$, \xout{de façon précise} \hand{(}tout comme les Z-\inserted{1-}gerbes peuvent être
interpétées comme des \xout{\unsure{1}}torseurs sous la Z-1-\over{\unsure{gerbe}}{champ} \inserted[typed]{de Picard strict,} des Z-torseurs, et sont
\xout{par conséquent} classifiées par le $H^2(X,Z)$\hand{)}. On voit déjà, bien sûr, poindre
ici l'oreille de la 3-catégorie de Picard stricte des 2-gerbes liées par Z,
ou (de façon équivalente) des 2-torseurs sous le 2-champ de \unsure{Picard} strict
des Z-1-gerbes; cette 3-catégorie de Picard stricte étant décrite par
$R\hand{\Gamma}_X(Z)$ tronqué en dimension 3, ayant comme invariants de cohomologie non
triviaux les $H^i(X,Z)$ $(0 \hand{\leqslant} i \hand{\leqslant} 3)$. Quant au 3-champ de Picard \unsure{correspondaht|correspondant},
il est décrit par une résolution injective de Z \hand{tronquée} en degré 3, alors
que le  2-champ de Picard précédent se décrivait en tronquant en degré 2.

2. Si M et N sont deux faisceaux abéliens sur X, les \uline{champs de Picard}
(NB: 1-champs !) \uline{d'invariants} M \uline{et} N forment eux-même une\inserted{2-}catégorie de
\margin{$\times$} Picard stricte, représentée  sans doute par le complexe RHom(X(M),N) \inserted{(23)}
tronqué en \xout{dimension} degré 2, dont les invariants de chomologie non trivi-%
aux sont donc \hand{``}le drôle de $\mathrm{Ext}^2$\hand{''} de ma lettre à Deligne, et les Honnètes
$\mathrm{Ext}^i(M,N)$  $(0 \hand{\leqslant} i \hand{\leqslant} 1)$. Les champs de Picard stricts forment une sous-2-caté%
gorie de Picard pleine, représentée par RHom(M,N) tronqué en degré 2,
d'invariants les $\mathrm{Ext}^i(M;N)$ $(0 \hand{\leqslant} i \hand{\leqslant} 2)$. Bien sûr, $\mathrm{Ext}^2$ donne les\inserted{0-}objets à
équivalence près, $\mathrm{Ext}^1$ les automorphismes à isomorphisme près de l'objet
nul, $\mathrm{Ext}^0$ les automorphismes de l'automorphisme identique audit ... Je n'ai
pas réfléchi à une bonne \unsure{interprétation} géométrique de la n-catégorie de
Picard associée à RHom(M,N) tronqué en degré n, et encore moins bien \xout{s}sûr
pour RHom(X(M),N), mais sans doute il faut regarder dans la direction des
n-champs de Picard.

3. Soit G un Groupe sur X, opérant sur un faisceau abélien N. Les \uline{champs}
\uline{en} (Gr)-\uline{catégories sur} X \uline{liés par} (G,N) forment une 2-catégorie de Picard,

\mspage{39}
